\section{Materials and methods}
\label{sec:methods}

\subsection{Instrument}
\label{sec:instrument}

The instrument under test is HunStat2, a single-channel three-electrode
potentiostat built on the Analog Devices AD5941 electrochemical analog front end
(AFE) and a Seeed Studio XIAO RP2040 microcontroller, on the
\texttt{STEI\_Bstat\_v1.2} carrier board. The AD5941 integrates two independent
excitation and sensing paths~\cite{devices2024ad5940}. The \emph{low-power} (LP)
loop comprises a 12-bit $V_{\mathrm{bias}}$ and 6-bit $V_{\mathrm{zero}}$
low-power DAC, a low-power potentiostat amplifier (LPPA) that drives the counter
electrode in feedback with the reference electrode, and a low-power
transimpedance amplifier (LPTIA) with a selectable internal feedback resistor
$\Rtia$ that senses working-electrode current. The \emph{high-speed} (HS) loop
provides sinusoidal excitation and a separate transimpedance stage for impedance
spectroscopy. Direct-current and staircase methods use the LP loop.

The board carries a bank of jumper-selectable calibration resistors
($\Rcal$) of \SI{200}{\ohm}, \SI{4.7}{\kilo\ohm} and \SI{10}{\kilo\ohm}, plus one
unpopulated position. These are not in the measurement signal path; they are the
reference against which the AFE measures its own transimpedance gain. On the
board used here, the \SI{200}{\ohm} position (\texttt{RCAL1}) is fitted.

Firmware is C/C++ built with \texttt{arduino-cli} 1.2.0 against the
\texttt{rp2040:rp2040} core 5.6.0, targeting
\texttt{rp2040:rp2040:seeed\_xiao\_rp2040}. Cyclic voltammetry executes through a
hardware-sequencer engine derived from the Analog Devices \texttt{AD5940\_Ramp}
reference example; chronoamperometry (CA), square-wave voltammetry (SWV),
differential pulse voltammetry (DPV), impedance spectroscopy (EIS) and
open-circuit potential (OCP) are implemented as direct-register C++ classes. The
host communicates over USB-CDC at \SI{1}{\mega\baud} with a compact ASCII command
protocol. All measurements reported here were driven from a Python
(\texttt{pyserial}) acquisition script rather than the graphical client, so that
every sweep is recorded with the settings and the calibrated transimpedance
value in effect at the time.

\subsection{Reference cell}
\label{sec:cell}

A precision passive network was used rather than an electrochemical cell.
Screen-printed electrodes vary between units and with reuse~\cite{anshori2025ganestat,
finvsgar2018reusability}, so an electrode-based comparison confounds electrode
variability with instrument performance; a passive network makes instrument
behaviour the dominant source of variation~\cite{kasper2022electrochemical} and,
critically for this work, has an absolute impedance known from component nominal
values rather than from a second instrument.

The external dummy cell exposes three branches from a shared counter/reference
node. The branch used throughout is \textbf{we-C}, a Randles topology
$R_s + (R_{ct} \parallel C_{dl})$ with
$R_s = \SI{560}{\ohm}$,
$R_{ct} = \SI{10}{\kilo\ohm}$ (\SI{0.1}{\percent} tolerance) and
$C_{dl} = \SI{33}{\nano\farad}$
(Fig.~\ref{fig:setup}). The remaining branches---a bare \SI{10}{\kilo\ohm}
\SI{0.1}{\percent} resistor (we-B) and an anti-parallel BAS16 diode pair in
parallel with \SI{33}{\nano\farad} in series with \SI{560}{\ohm}
(we-A)---were not connected.

For the direct-current methods considered here the relevant reference quantity
is the zero-frequency resistance,
\begin{equation}
R_{\mathrm{DC}} = R_s + R_{ct} = \Rnom,
\label{eq:rdc}
\end{equation}
whose uncertainty is dominated by $R_s$. Under a linear potential ramp of rate
$\nu$ the same network has the closed-form steady-state response
\begin{equation}
I(t) = \frac{\nu t}{R_s + R_{ct}}
     \;\pm\; \nu\,\frac{R_{ct}^{2} C_{dl}}{(R_s + R_{ct})^{2}},
\label{eq:ramp}
\end{equation}
the sign of the second term following the sweep direction. The forward and
reverse branches of a cyclic voltammogram are therefore predicted to be parallel
straight lines of slope $1/R_{\mathrm{DC}}$, separated by
$\Delta I = 2\nu C_{\mathrm{eff}}$ with
$C_{\mathrm{eff}} = R_{ct}^{2}C_{dl}/(R_s+R_{ct})^{2} = \SI{29.6}{\nano\farad}$.
Both the slope and the separation are thus predicted analytically from component
values, with no free parameters and no reference instrument.

\subsection{Where the current scale comes from}
\label{sec:traceability}

The quantity a potentiostat reports is not measured directly; it is the end of a
chain. For the LP loop, the reported current is
\begin{equation}
I_{\mathrm{rep}}
  = \frac{N_{\mathrm{ADC}}}{2^{15}}\cdot\frac{V_{\mathrm{ref}}}{G_{\mathrm{PGA}}\,\hat{R}_{\mathrm{TIA}}},
\label{eq:current}
\end{equation}
where $N_{\mathrm{ADC}}$ is the offset-binary ADC code referred to midscale,
$V_{\mathrm{ref}} = \SI{1.82}{\volt}$, $G_{\mathrm{PGA}}$ the programmable-gain
setting, and $\hat{R}_{\mathrm{TIA}}$ the \emph{measured} transimpedance gain.
The AD5941 does not assume its internal $\Rtia$; it measures it, by exciting a
known resistor and taking a ratio,
\begin{equation}
\hat{R}_{\mathrm{TIA}}
  = R_{\mathrm{CAL}}^{\,\mathrm{decl}}\cdot
    \frac{\mathrm{DFT}_{\mathrm{TIA}}}{\mathrm{DFT}_{\mathrm{CAL}}},
\label{eq:rtiacal}
\end{equation}
in which $R_{\mathrm{CAL}}^{\,\mathrm{decl}}$ is a \emph{software constant}
declaring the value of the physically fitted resistor. Combining
(\ref{eq:current}) and (\ref{eq:rtiacal}), if the declared value differs from the
fitted value by a factor $\gamma = R_{\mathrm{CAL}}^{\,\mathrm{decl}} /
R_{\mathrm{CAL}}^{\,\mathrm{fitted}}$, every reported current is scaled by
$1/\gamma$. This is the gain error $k$ of Section~\ref{sec:intro}: a single
mis-declared constant, entering multiplicatively, invisible to linearity.

A second-order consequence matters in practice. The excitation amplitude used
for the calibration in (\ref{eq:rtiacal}) is itself sized from the declared
value, as $V_{\mathrm{exc}} \approx 2000\times0.8\times
R_{\mathrm{CAL}}^{\,\mathrm{decl}}/\Rtia$~mV$_{pp}$. Over-declaring $\Rcal$
therefore also over-drives the calibration, and once the transimpedance output
saturates the ratio in (\ref{eq:rtiacal}) is no longer $\gamma$-proportional.
The resulting gain error is then neither predictable nor correctable after the
fact---it must be prevented.

\subsection{Validation protocol}
\label{sec:protocol}

We propose and apply the five checks below. They are ordered so that the checks
that constrain absolute scale come first, and none requires a reference
instrument. Acceptance windows were fixed before the measurements, adapting the
$\pm\SI{2}{\percent}$ resistive tolerance used by a commercial verification-cell
manual~\cite{gamry2011universal}.

\begin{description}
\item[V1 --- Scale traceability.] Confirm that every software constant that
  enters the gain chain matches the hardware actually fitted, and require the
  instrument to \emph{report} its calibrated $\hat{R}_{\mathrm{TIA}}$ so the
  scale is observable rather than assumed. Pass: declared $\Rcal$ equals the
  fitted resistor, and $\hat{R}_{\mathrm{TIA}}$ lies within the manufacturer's
  stated spread of its nominal setting.
\item[V2 --- Absolute accuracy.] Fit $I$ against $E$ and compare the recovered
  resistance $R = (\mathrm{d}I/\mathrm{d}E)^{-1}$ against the nominal value of
  Eq.~(\ref{eq:rdc}). Pass: $|R - R_{\mathrm{DC}}|/R_{\mathrm{DC}} \le
  \SI{2}{\percent}$. This is the check that linearity cannot substitute for.
\item[V3 --- Residual structure.] Report $R^2$ \emph{and} the root-mean-square
  fit residual expressed in units of the current least significant bit,
  $\mathrm{LSB} = V_{\mathrm{ref}}/(2^{15}G_{\mathrm{PGA}}\hat{R}_{\mathrm{TIA}})$.
  A residual at or below the quantisation floor indicates the fit is limited by
  the converter rather than by the analog chain; structure above it indicates a
  systematic defect that $R^2$ alone will under-report.
\item[V4 --- Repeatability.] Acquire $n \ge 7$ replicate sweeps under identical
  settings and report the relative standard deviation of the recovered
  resistance. Pass: $\%\mathrm{RSD} \le \SI{0.5}{\percent}$.
\item[V5 --- Parameter invariance.] Repeat across scan rate, potential window
  and potential step. The recovered resistance is a property of the cell and
  must not depend on the acquisition settings. Pass: the spread of $R$ across
  each series is within the V2 window.
\end{description}

\subsection{Acquisition and analysis}
\label{sec:analysis}

Unless stated otherwise, cyclic voltammograms were acquired from
$-200$ to $+\SI{200}{\milli\volt}$ in \SI{5}{\milli\volt} steps at
\SI{200}{\milli\volt\per\second}, one cycle, with the LP transimpedance set to
its \SI{1}{\kilo\ohm} nominal position. Ten replicates were recorded for the
accuracy and repeatability analysis; three replicates at each setting for the
$\Rcal$ mis-declaration series; and single sweeps across the scan-rate,
window and step series.

Each sweep was fitted by ordinary least squares over all points of both
branches. Recovered resistance is the reciprocal of the fitted slope, converted
to ohms. Loop separation was evaluated by interpolating forward and reverse
branches onto a common potential grid over their overlapping span and taking the
mean difference. Analysis used NumPy 2.3.4 on Python 3.13. Raw sweeps, the
acquisition scripts and the analysis that generates every figure and table are
in the project repository.

\begin{figure}[!t]
\centering
\includegraphics[width=0.92\linewidth]{fig1_setup}
\caption{Measurement chain and reference cell. (a) The instrument: AD5941 AFE
with its low-power loop (LPDAC, LPPA, LPTIA) driving the cell, and the
jumper-selected $\Rcal$ network used to calibrate the transimpedance gain ---
note that $\Rcal$ sits outside the measurement path but sets the scale of every
reported current through Eq.~(\ref{eq:rtiacal}). (b) The we-C Randles branch of
the external dummy cell used as the traceable reference.}
\label{fig:setup}
\end{figure}
